\documentclass[10pt,letterpaper]{article}
\usepackage{cornell-notes}

\title{Clause 08.03: View Methods}
\author{}
\date{}
\setDocTitle{Clause 08.03: View Methods}
\setDocSubtitle{ISO/IEC/IEEE 42010:2022}
\setDocOwner{ISO 42010 Cornell Notes}
\setDocDate{\today}
\setCornellCollection{ISO/IEC/IEEE 42010:2022 Cornell Notes}
\setCornellUnitType{Clause}
\setCornellUnitNumber{08.03}
\setCornellUnitTitle{View Methods}
\hypersetup{pdftitle={Clause 08.03: View Methods},pdfsubject={Cornell notes},pdfauthor={}}

\begin{document}
\maketitle
\begin{CornellOverviewBox}{Overview}
{\Large\bfseries\textcolor{CNavy}{Section 8.3: View methods}}\\[3pt]
{\small\textbf{Classification:} Normative}\\
{\small\textbf{Learning objective:} Identify the mandatory, recommended, and permitted provisions governing view methods.}
\end{CornellOverviewBox}
\vspace{3pt}

\begin{CornellNotesTable}
\CornellNoteRow{Purpose}{Defines when view methods are needed and what guidance they should provide for transforming model or non-model information into view components.}
\CornellNoteRow{Requirement}{\begin{itemize}[leftmargin=*,nosep,topsep=1pt]
\item View methods must be defined in the specifications of model kinds (see 8.2), viewpoints (see 8.1), ADLs (see 7.2), and ADFs (see 7.1) when it is necessary to provide guidance on how views are constructed or used.
\end{itemize}}
\CornellNoteRow{Recommendation}{\begin{itemize}[leftmargin=*,nosep,topsep=1pt]
\item If a model is used as an information source when creating a view, a view method should define how AD elements will be portrayed in the view component, how model data are transformed or translated for use in the view component and how relationships between information from different sources are portrayed in the view component.
\item If a “non-model” is used as an information source when creating a view, a view method should define how its contents are portrayed in the view component as AD elements, how the non-model-related data are transformed or translated for use in the view component, and how relationships between information from different sources are portrayed in the view component.
\end{itemize}}
\CornellNoteRow{Permission}{\begin{itemize}[leftmargin=*,nosep,topsep=1pt]
\item A specification of an architecture viewpoint may include one or more view methods.
\end{itemize}}
\CornellNoteRow{Related sections}{8.2, 8.1, 7.2, 7.1}
\CornellNoteRow{Conformance check}{\begin{itemize}[leftmargin=*,nosep,topsep=1pt]
\item View methods must be defined in the specifications of model kinds (see 8.2), viewpoints (see 8.1), ADLs (see 7.2), and ADFs (see 7.1) when it is necessary to provide guidance on how views are constructed or used
\end{itemize}}
\CornellNoteRow{Active recall}{\begin{itemize}[leftmargin=*,nosep,topsep=1pt]
\item Which provisions in View methods are mandatory for conformance?
\item Which provisions are recommendations or permissions rather than requirements?
\end{itemize}}
\end{CornellNotesTable}


\begin{CornellSummaryBox}{Summary}
View methods are required in supporting specifications when guidance is necessary for constructing or using views. For model and non-model sources, recommended guidance covers how AD elements are portrayed, how source information is transformed, and how relationships among sources appear in view components. A viewpoint specification may include one or more view methods.
\end{CornellSummaryBox}

\end{document}
